\section{Gröbner Bases}

Let's move from ideals in $k[x]$ to ideals in $k[x_1, \dots, x_n]$. We are interested in qualitative and algorithmic theory of such ideals.
Our intention is to use this understanding for solving various natural problems from algebra and geometry (related to the varieties in $k^n$). 
We are inpired by the case $n =1$, we have analyzed above (and more interesting). After all algebra of $k[x_1, \dots, x_n]$ is richer than algebra of $k[x]$ 
and, respectively, the geometry of $k^n$ is richter than the geometry of $k^1$.